\section{Introducción}
El análisis cinemático permite relacionar el movimiento de las articulaciones de un robot con la posición y la velocidad de su efector final. Constituye la base para estudiar configuraciones alcanzables, singularidades y estrategias de seguimiento de trayectorias.

El presente trabajo desarrolla los contenidos del parcial de Robótica y Sistemas Autónomos a partir del manipulador antropomórfico de tres grados de libertad descrito por \textcite{ramzy2024}. Se abordarán la cinemática directa, inversa y diferencial, el análisis de singularidades y el control cinemático, acompañados de implementación en Python y simulación \parencite{consigna2026}.

El objetivo es formular y comprobar el modelo cinemático mediante casos reproducibles. El documento presenta el artículo de referencia, desarrolla cada tema y reúne resultados y conclusiones.
